\documentclass[a4paper,12pt]{article}
\usepackage[T2A]{fontenc}
\usepackage[utf8]{inputenc}
\usepackage[russian]{babel}
\usepackage{amsmath,amssymb,mathtools}
\usepackage{PTSans}
\renewcommand{\familydefault}{\sfdefault}
\usepackage{icomma}
\usepackage[a4paper,left=19mm,right=19mm,top=19mm,bottom=20mm]{geometry}
\usepackage{microtype}
\usepackage{tikz}
\usetikzlibrary{calc,arrows.meta,positioning,shapes.geometric}
\usepackage{tkz-euclide}
\usepackage{booktabs,tabularx,array}
\usepackage{enumitem}
\usepackage{fancyhdr}
\usepackage{setspace}
\usepackage{xcolor}
\usepackage{hyperref}
\usepackage{parskip}
\definecolor{accent}{HTML}{176B73}
\definecolor{darkaccent}{HTML}{124A50}
\definecolor{softgray}{HTML}{69767A}
\definecolor{textgray}{HTML}{30383B}
\definecolor{warm}{HTML}{B66A3C}
\hypersetup{colorlinks=true,linkcolor=darkaccent,urlcolor=darkaccent}
\setlength{\parindent}{0pt}
\setlength{\parskip}{5pt}
\setstretch{1.28}
\setlist[itemize]{leftmargin=6mm,itemsep=2pt,topsep=3pt}
\setlist[enumerate]{leftmargin=8mm,itemsep=4pt,topsep=3pt}
\pagestyle{fancy}
\fancyhf{}
\fancyfoot[C]{\small\color{softgray}\thepage}
\renewcommand{\headrulewidth}{0pt}
\newcommand{\sect}[1]{\vspace{2mm}{\Large\bfseries\color{darkaccent}#1}\par\vspace{1mm}{\color{accent!55}\rule{\linewidth}{0.55pt}}\par\vspace{1mm}}
\newcommand{\idea}[1]{\textbf{\color{darkaccent}#1}}
\begin{document}
\thispagestyle{empty}
\begin{center}
\vspace*{30mm}
{\small\color{softgray}МАТЕМАТИКА \textbullet{} ОГЭ \textbullet{} ГЕОМЕТРИЯ}\par
\vspace{11mm}
{\fontsize{28}{32}\selectfont\bfseries\color{darkaccent}Чек-лист}\par
\vspace{6mm}
{\LARGE\bfseries Высоты, касательные и отношения отрезков}\par
\vspace{7mm}
{\large Площади треугольников \textbullet{} прямоугольные треугольники \\ Теорема Фалеса и подобие}\par
\vspace{23mm}
{\large София Кальней}\par
\vspace{4mm}
{\normalsize 9 октября 2026 г.}\par
\vfill
{\small\color{textgray}Лёвин Артём Александрович эксклюзивно для Софии Кальней}\par
\vspace{20mm}
\end{center}
\begin{tikzpicture}[remember picture,overlay]
\draw[accent!55,line width=1.15pt]
 ([xshift=8mm,yshift=7mm]current page.south west)
 .. controls +(47mm,21mm) and +(-56mm,-5mm) ..
 ([xshift=-8mm,yshift=14mm]current page.south east);
\draw[warm!45,line width=0.7pt]
 ([xshift=16mm,yshift=5mm]current page.south west)
 .. controls +(55mm,17mm) and +(-66mm,-8mm) ..
 ([xshift=-16mm,yshift=11mm]current page.south east);
\end{tikzpicture}
\clearpage

\sect{1. Высота, расстояние до прямой и площадь}
\idea{Проверьте определения.} Расстояние от точки до прямой --- длина перпендикуляра, проведённого из этой точки к прямой. Высота треугольника --- такой перпендикуляр из вершины к прямой, содержащей противоположную сторону.

\[
S_{\triangle}=\frac12 ah,\qquad
p=\frac{a+b+c}{2},\qquad
S_{\triangle}=\sqrt{p(p-a)(p-b)(p-c)}.
\]
\textbf{Формула Герона} нужна, когда известны три стороны. Перед применением проверьте существование треугольника (сумма любых двух сторон больше третьей).

\idea{Составление площади из частей.} Если $M$ лежит \emph{внутри} $\triangle ABC$, проведите $AM$, $BM$, $CM$. Получится разбиение без наложений:
\[
S_{ABC}=S_{ABM}+S_{BCM}+S_{CAM}.
\]
Если $d_{AB},d_{BC},d_{CA}$ --- расстояния от $M$ до прямых $AB,BC,CA$, то
\[
\boxed{\;2S_{ABC}=AB\cdot d_{AB}+BC\cdot d_{BC}+CA\cdot d_{CA}.\;}
\]
Эта формула справедлива при $M$ внутри треугольника: все три площади складываются.

\textbf{Пример с занятия.} $AB=10$, $BC=17$, $AC=21$; расстояния от $M$ до $AC$ и $BC$ равны $2$ и $4$. Найдите расстояние $x$ до $AB$.
\[
p=24,\quad S=\sqrt{24\cdot14\cdot7\cdot3}=84;
\qquad10x+17\cdot4+21\cdot2=168.
\]
\[
10x=58\quad\Longrightarrow\quad\boxed{x=\frac{29}{5}}.
\]
\idea{Типичная ошибка.} Не подменяйте расстояние до стороны наклонным отрезком. Основание перпендикуляра к \emph{прямой} может находиться на продолжении стороны.

\begin{center}
\begin{tikzpicture}[scale=.73]
\tkzDefPoint(0,3.2){A}
\tkzDefPoint(-3.4,0){B}
\tkzDefPoint(3.6,0){C}
\tkzDefPoint(0.4,1.0){M}
\tkzDrawPolygon[color=textgray,line width=.9pt](A,B,C)
\tkzDrawSegments[color=accent,line width=.9pt](M,A M,B M,C)
\tkzDrawPoints[size=2,fill=darkaccent](A,B,C,M)
\tkzLabelPoint[above](A){$A$}
\tkzLabelPoint[below left](B){$B$}
\tkzLabelPoint[below right](C){$C$}
\tkzLabelPoint[right](M){$M$}
\end{tikzpicture}
\quad\begin{minipage}[c]{.47\linewidth}
\small \textbf{Три части одного треугольника.} Отрезки $AM$, $BM$, $CM$ превращают неизвестное расстояние от $M$ до стороны в высоту одного из трёх новых треугольников.
\end{minipage}
\end{center}
\clearpage
\sect{2. Равнобедренный треугольник: два способа найти площадь}
Пусть боковые стороны равны $b$, основание равно $a$, высота к основанию равна $h$, высота к боковой стороне --- $k$. Тогда
\[
\frac12 ah=\frac12 bk\quad\Longrightarrow\quad ah=bk.
\]
Высота к основанию является медианой, поэтому по теореме Пифагора
\[
\boxed{b^2=h^2+\left(\frac a2\right)^2.}
\]
\textbf{Пример с занятия.} $h=10$, $k=12$. Из равенства площадей $b=\frac56 a$. Тогда
\[
\frac{25a^2}{36}=100+\frac{a^2}{4}
\ \Longrightarrow\ a=15,\quad b=\frac{25}{2},\quad \boxed{S=75}.
\]
\idea{Приём.} Одна площадь, записанная через два разных основания и соответствующие высоты, даёт полезное отношение сторон.

\sect{3. Две касательные из одной точки}
Пусть из внешней точки $A$ к окружности с центром $O$ проведены касательные $AB$ и $AC$, $B,C$ --- точки касания, $F=AO\cap BC$.

\begin{minipage}[c]{0.54\linewidth}
\begin{itemize}
\item $OB\perp AB$ и $OC\perp AC$ (радиусы к касательным).
\item $AB=AC$ (касательные из точки $A$), $OB=OC$ (радиусы).
\item Точки $A,O$ лежат на серединном перпендикуляре к $BC$; значит $AO\perp BC$ и $BF=FC$.
\end{itemize}
\end{minipage}\hfill
\begin{minipage}[c]{0.42\linewidth}
\centering
\begin{tikzpicture}[scale=0.72]
\tkzDefPoint(0,0){O}
\tkzDefPoint(4.25,0){A}
\pgfmathsetmacro{\yy}{sqrt(2.89-0.68*0.68)}
\tkzDefPoint(0.68,\yy){B}
\tkzDefPoint(0.68,-\yy){C}
\tkzDefPoint(0.68,0){F}
\tkzDrawCircle[color=accent,line width=.8pt](O,B)
\tkzDrawSegments[line width=.8pt](A,B A,C O,B O,C A,O B,C)
\tkzDrawPoints[fill=darkaccent,size=2.3](A,B,C,O,F)
\tkzLabelPoint[right](A){$A$}
\tkzLabelPoint[above](B){$B$}
\tkzLabelPoint[below](C){$C$}
\tkzLabelPoint[left](O){$O$}
\tkzLabelPoint[below right](F){$F$}
\tkzMarkRightAngle[size=.18](A,B,O)
\tkzMarkRightAngle[size=.18](O,C,A)
\tkzMarkRightAngle[size=.18](B,F,O)
\end{tikzpicture}
\end{minipage}

\textbf{Пример с занятия.} $OB=7$, $AO=25$. В $\triangle AOB$:
\[
AB=\sqrt{AO^2-OB^2}=\sqrt{625-49}=24.
\]
Поскольку $BF$ --- высота к гипотенузе $AO$, найдём её, приравняв два выражения для площади $\triangle AOB$:
\[
\frac12 AB\cdot OB=\frac12 AO\cdot BF,
\qquad BF=\frac{24\cdot7}{25}=\frac{168}{25}.
\]
\[
\boxed{BC=2BF=\frac{336}{25}.}
\]
\idea{Типичная ошибка.} Перпендикулярность $AO$ и $BC$ сначала \emph{обоснуйте}, затем используйте $BF=FC$.

\sect{4. Теорема Фалеса и отношения отрезков}
Если две параллельные прямые пересекают стороны угла, они отсекают \textbf{пропорциональные отрезки}. В треугольнике $ABC$ при $D\in AB$, $E\in AC$ и $DE\parallel BC$:
\[
\triangle ADE\sim\triangle ABC,
\qquad\frac{AD}{AB}=\frac{AE}{AC}=\frac{DE}{BC},
\qquad\boxed{\frac{AD}{DB}=\frac{AE}{EC}}.
\]
\idea{Проверяйте соответствие.} Сопоставляйте отрезки, занимающие одинаковое положение на сторонах угла.

\sect{5. Медиана и параллельная прямая}
Пусть $AM$ --- медиана $\triangle ABC$, $K\in AM$, а прямая через $K$, параллельная $AC$, пересекает $MC$ в точке $J$.

\begin{minipage}[c]{0.55\linewidth}
\begin{enumerate}
\item Из $AM$ --- медиана следует $BM=MC$.
\item Из $KJ\parallel AC$ следует $\triangle MKJ\sim\triangle MAC$.
\item Поэтому $\displaystyle\frac{KM}{KA}=\frac{MJ}{JC}$.
\item Выразите отрезки через части и сложите смежные отрезки на $BC$.
\end{enumerate}
\end{minipage}\hfill
\begin{minipage}[c]{0.43\linewidth}
\centering
\begin{tikzpicture}[scale=.93]
\tkzDefPoint(0.3,3.5){A}
\tkzDefPoint(-2.6,0){B}
\tkzDefPoint(3.2,0){C}
\tkzDefPoint(0.3,0){M}
\tkzDefPoint(0.3,2.625){K}
\tkzDefPoint(2.475,0){J}
\tkzDrawPolygon[line width=.8pt](A,B,C)
\tkzDrawSegments[line width=.8pt](A,M)
\tkzDrawSegments[color=accent,line width=1.4pt](K,J)
\tkzDrawPoints[fill=darkaccent,size=2.5](A,B,C,M,K,J)
\tkzLabelPoint[above](A){$A$}
\tkzLabelPoint[left](B){$B$}
\tkzLabelPoint[right](C){$C$}
\tkzLabelPoint[below](M){$M$}
\tkzLabelPoint[left](K){$K$}
\tkzLabelPoint[below](J){$J$}
\end{tikzpicture}
\end{minipage}

\textbf{Пример с занятия.} Если $AK:KM=1:3$, то из пропорциональности $MJ:JC=3:1$. Обозначим $MJ=3t$, $JC=t$. Тогда $MC=4t$ и $BM=4t$.
\[
BJ=BM+MJ=7t,\quad\boxed{BJ:JC=7:1}.
\]
\idea{Обобщение.} Для $AK:KM=p:q$ (положительные $p,q$):
\[
MJ:JC=q:p,\qquad\boxed{BJ:JC=(p+2q):p}.
\]

\sect{6. Как действовать в задаче}
\begin{enumerate}[itemsep=1pt]
\item Отметьте известные длины, перпендикуляры и параллели.
\item Найдите полезные треугольники; при необходимости проведите дополнительные отрезки.
\item Выберите формулу: площадь, Пифагор, равенство или подобие треугольников.
\item Введите неизвестную или доли отрезков, составьте равенство и проверьте ответ.
\end{enumerate}
\clearpage
\thispagestyle{plain}
\begin{center}{\Large\bfseries\color{darkaccent}Алгоритм: какой метод выбрать?}\end{center}
\vspace{8mm}
\begin{center}
\begin{tikzpicture}[
 >={Latex[length=2.6mm]},
 node distance=11mm,
 begin/.style={ellipse,draw=darkaccent,thick,align=center,minimum width=48mm,minimum height=10mm,fill=accent!5},
 question/.style={diamond,draw=accent,thick,aspect=2.5,inner sep=2pt,align=center,text width=46mm,font=\small},
 action/.style={rectangle,rounded corners=2pt,draw=accent,thick,align=center,text width=48mm,minimum height=13mm,inner sep=4pt,font=\small},
 arrow/.style={->,thick,color=textgray}
]
\node[begin] (start) {Чертёж и данные};
\node[question,below=13mm of start] (par) {Есть параллельные\\прямые?};
\node[action,right=18mm of par] (paract) {Найдите подобные треугольники,\\составьте отношения};
\node[question,below=18mm of par] (hei) {Есть высоты или\\расстояния?};
\node[action,right=18mm of hei] (heiact) {Используйте площади,\\Пифагора или Герона};
\node[question,below=18mm of hei] (tan) {Есть касательные\\к окружности?};
\node[action,right=18mm of tan] (tanact) {Проведите радиусы к точкам касания, найдите прямые углы};
\node[action,below=17mm of tan] (other) {Ищите полезные дополнительные треугольники};
\node[begin,below=12mm of other] (end) {Составьте равенство\\и проверьте результат};
\draw[arrow] (start)--(par);
\draw[arrow] (par.east)--node[above,font=\scriptsize]{да}(paract.west);
\draw[arrow] (par.south)--node[right,font=\scriptsize]{нет}(hei.north);
\draw[arrow] (hei.east)--node[above,font=\scriptsize]{да}(heiact.west);
\draw[arrow] (hei.south)--node[right,font=\scriptsize]{нет}(tan.north);
\draw[arrow] (tan.east)--node[above,font=\scriptsize]{да}(tanact.west);
\draw[arrow] (tan.south)--node[right,font=\scriptsize]{нет}(other.north);
\draw[arrow] (other)--(end);
\coordinate (merge) at ($(end.south)+(0,-5mm)$);
\draw[arrow] (paract.east)--++(5mm,0)|-($(end.east)+(0,1mm)$);
\draw[arrow] (heiact.east)--++(9mm,0)|-($(end.east)+(0,1mm)$);
\draw[arrow] (tanact.east)--++(13mm,0)|-($(end.east)+(0,1mm)$);
\end{tikzpicture}
\end{center}
\clearpage

\sect{Тренировочный блок. Путь А --- 10 заданий}
\textit{Выполняйте задания последовательно. Записывайте обоснования построений и используемых теорем. Решения в этом блоке не приведены.}

\begin{enumerate}[itemsep=7pt]
\item В прямоугольном треугольнике $ABC$ стороны $AB=12$, $AC=5$, $BC=13$. Точка $M$ лежит внутри треугольника. Расстояния от $M$ до прямых $AB$ и $AC$ равны $2$ и $3$. Найдите расстояние до прямой $BC$.
\item В треугольнике $ABC$ даны $AB=13$, $BC=14$, $AC=15$. Точка $M$ лежит внутри треугольника; расстояния от $M$ до $AC$ и $BC$ равны $2$ и $3$. Найдите расстояние от $M$ до $AB$.
\item В равнобедренном треугольнике высота к основанию равна $20$, а высота к боковой стороне равна $24$. Найдите площадь треугольника.
\item Из точки $A$ к окружности радиуса $8$ проведена касательная $AB$ ($B$ --- точка касания). Если расстояние от $A$ до центра $O$ равно $17$, найдите $AB$.
\item Из точки $A$ к окружности радиуса $5$ проведены касательные $AB$ и $AC$, причём $AO=13$. Найдите расстояние между точками касания $B$ и $C$.
\item На сторонах угла с вершиной $P$ отмечены точки $A,B$ и $C,D$ соответственно; на лучах они расположены в порядке $P,A,B$ и $P,C,D$. Известно, что $AC\parallel BD$, $PA=4$, $AB=6$, $PC=7$. Найдите $CD$.
\item В треугольнике $ABC$ точки $D\in AB$ и $E\in AC$ таковы, что $DE\parallel BC$. Если $AD=6$, $DB=9$, $AE=4$, найдите $EC$.
\item Из внешней точки $A$ к окружности с центром $O$ проведены касательные $AB$ и $AC$ ($B,C$ --- точки касания). Докажите, что прямая $AO$ перпендикулярна $BC$.
\item В треугольнике $ABC$ проведена медиана $AM$. Точка $K$ лежит на $AM$ и $AK:KM=2:3$. Через $K$ проведена прямая, параллельная $AC$, пересекающая $MC$ в точке $J$. Найдите $BJ:JC$.
\item В треугольнике $ABC$ медиана $AM$ содержит точку $K$, причём $AK:KM=3:2$. Прямая через $K$, параллельная $AC$, пересекает $MC$ в точке $J$. Если $BJ=21$, найдите $JC$.
\end{enumerate}
\clearpage

\sect{Ответы и опорные шаги}
\renewcommand{\arraystretch}{1.4}
\begin{tabularx}{\linewidth}{@{}c p{27mm} X@{}}
\toprule
\textbf{№} & \textbf{Ответ} & \textbf{Ключевая идея для проверки} \\
\midrule
1 & $\dfrac{21}{13}$ & $S=30$; $12\cdot2+5\cdot3+13x=60$. \\
2 & $\dfrac{96}{13}$ & $p=21$, $S=84$; $13x+14\cdot3+15\cdot2=168$. \\
3 & $300$ & Если основание $a$, боковая сторона $b$, то $20a=24b$ и $b^2=20^2+(a/2)^2$; $a=30$. \\
4 & $15$ & Примените Пифагора в прямоугольном $\triangle AOB$. \\
5 & $\dfrac{120}{13}$ & $AB=12$, $BF=\dfrac{AB\cdot OB}{AO}=\dfrac{60}{13}$, $BC=2BF$. \\
6 & $\dfrac{21}{2}$ & $\dfrac{PA}{AB}=\dfrac{PC}{CD}$. \\
7 & $6$ & $\dfrac{AD}{DB}=\dfrac{AE}{EC}$. \\
8 & $AO\perp BC$ & $AB=AC$ и $OB=OC$; точки $A,O$ лежат на серединном перпендикуляре к $BC$. \\
9 & $4:1$ & $MJ:JC=3:2$, $BM=MC$; $BJ:JC=8:2$. \\
10 & $9$ & $MJ:JC=2:3$, $BM=MC$; $BJ:JC=7:3$. \\
\bottomrule
\end{tabularx}
\vspace{8mm}
\idea{Итог самопроверки.} Вы можете объяснить каждое дополнительное построение, указать основание для подобия или перпендикулярности и получить ответ без опоры на рисунок «на глаз».
\end{document}
